
\documentclass[preprint,12pt]{elsarticle}
\usepackage{multirow}
\usepackage[table]{xcolor} 
\usepackage[normalem]{ulem}
\usepackage{colortbl}
\usepackage{amsmath,amssymb,amsfonts}
\usepackage{algpseudocode}
\usepackage{graphicx}
\usepackage{textcomp}
\usepackage{booktabs}
\usepackage{bm}
\usepackage{algorithm}
\newcommand{\abcell}[2]{\shortstack[r]{#1\\[-1pt]{\scriptsize(#2)}}}
\makeatletter
\AtBeginDocument{\DeclareMathVersion{bold}
\SetSymbolFont{operators}{bold}{T1}{times}{b}{n}
\DeclareSymbolFont{NewLetters}{T1}{cmr}{m}{it}
\SetSymbolFont{NewLetters}{bold}{T1}{times}{b}{it}
\SetMathAlphabet{\mathrm}{bold}{T1}{times}{b}{n}
\SetMathAlphabet{\mathit}{bold}{T1}{times}{b}{it}
\SetMathAlphabet{\mathbf}{bold}{T1}{times}{b}{n}
\SetMathAlphabet{\mathtt}{bold}{OT1}{pcr}{b}{n}
\SetSymbolFont{symbols}{bold}{OMS}{cmsy}{b}{n}
\renewcommand\boldmath{\@nomath\boldmath\mathversion{bold}}}
\newtheorem{proposition}{Proposition}
\newtheorem{lemma}{Lemma}
\newtheorem{remark}{Remark}
\newtheorem{proof}{Proof}
\newtheorem{corollary}{Corollary}
\newcommand{\OurPCC}{0.XXX}
\newcommand{\OurRMSE}{XX.XX} % in %
\newcommand{\OurMAE}{XX.XX}  % in %
\newcommand{\OurLatency}{XX.X} % ms per update in streaming mode (step-wise)

\journal{Knowledge-Based Systems}

\begin{document}

\begin{frontmatter}

%% Title, authors and addresses

%% use the tnoteref command within \title for footnotes;
%% use the tnotetext command for theassociated footnote;
%% use the fnref command within \author or \affiliation for footnotes;
%% use the fntext command for theassociated footnote;
%% use the corref command within \author for corresponding author footnotes;
%% use the cortext command for theassociated footnote;
%% use the ead command for the email address,
%% and the form \ead[url] for the home page:
%% \title{Title\tnoteref{label1}}
%% \tnotetext[label1]{}
%% \author{Name\corref{cor1}\fnref{label2}}
%% \ead{email address}
%% \ead[url]{home page}
%% \fntext[label2]{}
%% \cortext[cor1]{}
%% \affiliation{organization={},
%%             addressline={},
%%             city={},
%%             postcode={},
%%             state={},
%%             country={}}
%% \fntext[label3]{}

\title{From Electrode Geometry to Adaptive Coupling: HASTE for Streaming EEG Fatigue Monitoring}


\author{Qiyu Niu}
\affiliation{organization={Academy for Engineering and Technology, Fudan University},
            city={Shanghai},
            postcode={200433},
            country={China}}

\affiliation{organization={Institute of Science and Technology for Brain-Inspired Intelligence, Fudan University},
            city={Shanghai},
            postcode={200433},
            country={China}}

\author{Yi Xie}
\affiliation{organization={Department of Electrical \& Computer Engineering, University of Arizona},
            city={Tucson},
            postcode={85721},
            state={Arizona},
            country={USA}}

\author{Siao Liu}
\affiliation{organization={Future Science and Engineering College, Soochow University},
            city={Suzhou},
            postcode={215006},
            state={Jiangsu},
            country={China}}

\author{Shouyan Wang*}
\affiliation{organization={Academy for Engineering and Technology, Fudan University},
            city={Shanghai},
            postcode={200433},
            country={China}}

\affiliation{organization={Institute of Science and Technology for Brain-Inspired Intelligence, Fudan University},
            city={Shanghai},
            postcode={200433},
            country={China}}

\author{Yingnan Nie*}
\affiliation{organization={Institute of Science and Technology for Brain-Inspired Intelligence, Fudan University},
            city={Shanghai},
            postcode={200433},
            country={China}}

\affiliation{organization={School of Rehabilitation Sciences and Engineering, University of Health and Rehabilitation Sciences},
            city={Qingdao},
            postcode={266113},
            state={Shandong},
            country={China}}

\affiliation{country={*Corresponding authors}}            

\begin{abstract}
Streaming EEG fatigue monitoring is challenging because fatigue signatures are both frequency-dependent and spatially distributed across electrodes, while deployment requires strictly causal updates with bounded latency and memory. Existing Transformers weakly exploit electrode geometry and incur history-dependent cost, whereas graph models typically rely on fixed or unconstrained connectivity that is difficult to adapt as cognitive state changes. We propose HASTE, a hybrid-attention streaming topology-aware estimator that separates where interactions are allowed from how strongly they adapt: a physical $k$NN graph constrains spatial attention to anatomically feasible neighborhoods, and an online band-aware functional prior biases within-neighborhood attention through an additive log-bias. Causal spectral fusion tracks time-varying band relevance, and causal local temporal attention with a fixed look-back window enables constant-memory inference with respect to session length. Under leave-one-subject-out streaming evaluation on SEED-VIG, SADT, and MPD-DF, HASTE consistently outperforms six baselines, achieving PCCs of 0.827, 0.751, and 0.714, respectively, with ${\sim}$2.8\,ms per update and ${\sim}$47\,MB GPU memory on a single GPU.
\end{abstract}

\begin{highlights}
\item A physical $k$NN graph constrains spatial attention to anatomically feasible neighborhoods.
\item Band-aware functional priors bias spatial attention through a tunable additive log-bias.
\item Causal spectral fusion tracks time-varying band relevance and mixes band-specific priors online.
\item Causal local temporal attention with a fixed look-back window enables constant-memory inference.
\item HASTE improves LOSO streaming fatigue regression on SEED-VIG, SADT, and MPD-DF.
\end{highlights}

\begin{keyword}
EEG \sep fatigue monitoring \sep streaming inference \sep
topology-aware attention
\end{keyword}
\end{frontmatter}


\section{Introduction}
\label{sec:intro}

Fatigue monitoring is critical in safety-sensitive settings such as prolonged driving and industrial operation, where performance degradation can accumulate before it becomes apparent \cite{Williamson2011FatigueSafety,Sikander2019DriverFatigueReview}. 
A deployable monitor must therefore update strictly causally in real time, with bounded per-step latency and memory that does not grow with session length.

EEG is well suited to this problem because fatigue-related activity exhibits two coupled structures: frequency-dependent rhythms and state-varying inter-electrode coupling \cite{Lal2001DriverFatigue,Eoh2005DrowsinessEEG}. 
Continuous supervision is often available through surrogate measures such as PERCLOS or reaction-time-derived fatigue indices \cite{Dinges1998PERCLOS,Zheng2017CriticalBands,cao2019multi,Li2026MPDDF}. 
We consider the subject-independent streaming regression setting: at time step $t$, the model receives the current window of multi-band EEG features from $C$ electrodes and predicts a continuous fatigue score using only observations up to $t$. 
Throughout, causal refers to the model update given window-level features; the same requirement can be imposed on feature extraction in deployment.

However, existing architectures do not fit this setting well. 
First, Transformers usually model EEG as a generic token sequence and often weakly exploit electrode geometry. 
Second, their temporal cost grows with history length unless additional constraints are imposed \cite{vaswani2017attention,Beltagy2020Longformer}. 
Third, graph-based models incorporate spatial structure, but they typically rely on fixed connectivity or dense learned interactions, making state-varying coupling hard to adapt and hard to control \cite{kipf2016semi,velivckovic2017graph}. 
The central question is therefore whether a streaming EEG model can remain band-aware and state-adaptive while respecting electrode geometry and bounded deployment cost.

To answer this question, we propose \textit{HASTE}, a \textit{H}ybrid-\textit{A}ttention \textit{S}treaming \textit{T}opology-aware \textit{E}stimator. 
HASTE first builds a sparse physical graph from electrode locations and restricts spatial attention to a $k$NN neighborhood, so information exchange remains anatomically plausible \cite{Klem1999TenTwenty}. 
Within this feasible set, it computes a row-stochastic functional coupling prior online and injects it into the attention logits as an additive log-bias, making the effect of prior knowledge explicit and tunable. 
To follow time-varying spectral relevance, HASTE includes a causal spectral fusion module that maintains a recurrent state and mixes band-specific priors online \cite{Cho2014GRU}. 
For temporal modeling, it uses causal local attention with a fixed look-back window and a key--value cache, enabling bounded-latency, constant-memory inference with respect to session length \cite{dai2019transformer,Beltagy2020Longformer}. 
Figure\ref{figure1} summarizes the streaming setting, the main HASTE modules, and the interpretable outputs exposed by the model, including the band weights $\boldsymbol{\alpha}_t$, the prior-mixing gate $g_t$, and stage-dependent spatial salience patterns.

\begin{figure*}[!t]
\centering
\label{figure1}
\includegraphics[width=0.98\linewidth]{Figure1.pdf}
\caption{High-level overview of HASTE in the streaming EEG fatigue monitoring setting. (A) Each update uses the current multi-band EEG window together with the electrode montage. (B) The model combines causal spectral fusion, knowledge-guided spatial attention, and causal local temporal attention. (C) HASTE outputs a fatigue estimate and exposes interpretable signals, including band weights $\boldsymbol{\alpha}_t$, the prior-mixing gate $g_t$, and stage-dependent spatial salience patterns.}
\end{figure*}

We evaluate HASTE under a unified leave-one-subject-out streaming pipeline on SEED-VIG, SADT, and MPD-DF \cite{Zheng2017CriticalBands,cao2019multi,Li2026MPDDF}. 
Beyond regression accuracy, we profile per-update latency and peak GPU memory as session length grows, and perform ablations and interpretability analyses to isolate the contribution of each mechanism. 
Across all three datasets, HASTE consistently improves over strong baselines while retaining millisecond-level update time and constant-memory behavior with respect to session length.

In summary, our contributions are fourfold:
\begin{itemize}
\item We formulate streaming EEG spatial modeling for subject-independent fatigue regression as a constrained problem, and show that topology-constrained stochastic priors with additive log-bias attention provide explicit prior control and stability under prior drift.

\item Based on this formulation, we develop a knowledge-guided spatial design that combines a physical $k$NN graph, band-aware functional priors, and anatomical--functional prior fusion for adaptive yet anatomically feasible coupling.

\item We propose HASTE, which instantiates this design in a hybrid-attention spatio-temporal architecture with causal spectral fusion, prior-biased spatial attention, and bounded causal local temporal attention for constant-memory streaming inference.

\item We establish a unified LOSO streaming evaluation on SEED-VIG, SADT, and MPD-DF, showing that HASTE consistently outperforms strong baselines while retaining predictable deployment cost.
\end{itemize}

\section{Related Work}
\label{sec:related_work}

Our work is most closely related to three research directions: continuous EEG-based fatigue and vigilance estimation, topology-aware spatio-temporal EEG modeling, and efficient causal sequence modeling for streaming inference.

\subsection{Continuous EEG-Based Fatigue and Vigilance Estimation}

Fatigue monitoring has long been studied in safety-critical scenarios such as driving and transportation, and prior surveys summarize sensing modalities, fatigue proxies, and algorithmic pipelines \cite{Sikander2019DriverFatigueReview,Williamson2011FatigueSafety}. 
Among physiological signals, EEG is particularly attractive because fatigue-related changes are reflected in vigilance-related spectral activity and cortical dynamics \cite{Lal2001DriverFatigue,Eoh2005DrowsinessEEG}. 
A practical difficulty is that fatigue evolves continuously rather than switching cleanly between a few discrete states. 
Accordingly, continuous supervision has been introduced through ocular or behavioral proxies. 
PERCLOS is one of the most established continuous alertness measures \cite{Dinges1998PERCLOS}, and early multimodal vigilance estimation work already formulated EEG-based monitoring as a continuous regression problem rather than a pure classification task \cite{Zheng2017CriticalBands}. 
Public benchmarks reinforce this setting: SEED-VIG provides continuous vigilance targets, SADT offers sustained-attention driving EEG with reaction-time events that can be mapped to continuous indices, and MPD-DF further extends subject scale, signal modalities, and physician-annotated fatigue labels \cite{Zheng2017CriticalBands,cao2019multi,Li2026MPDDF}. 
Recent work has also begun to explore cross-dataset fatigue regression \cite{Yuan2023CrossDatasetFatigue}. 
Nevertheless, explicit streaming requirements---strict causality, bounded per-step latency, and memory that does not grow with session length---remain secondary in most of this literature.

\subsection{Topology-Aware Spatio-Temporal EEG Modeling}

Deep EEG representation learning has progressed from compact and deep convolutional models to recurrent sequential models \cite{Roy2019DeepLearningEEGReview,Lawhern2018EEGNet,Schirrmeister2017DeepConvNets,Hochreiter1997LSTM}. 
More recently, Transformer-style EEG architectures have been explored to capture broader temporal dependencies and global feature interactions \cite{vaswani2017attention,song2023eeg}. 
These studies show the value of attention for EEG modeling, but spatial relations are typically handled through convolutions or unrestricted attention rather than an explicit topology-constrained interaction graph. 

In parallel, graph neural networks provide a natural framework for topology-aware EEG modeling. 
Foundational graph models such as GCN and GAT establish neighborhood aggregation and attention-based message passing on graphs \cite{kipf2016semi,velivckovic2017graph}, and EEG-specific dynamic graph constructions such as DGCNN demonstrate the value of adaptive connectivity \cite{Song2018DGCNN}. 
This literature reveals a core trade-off: fixed topology promotes anatomical plausibility and stability, whereas fully learned connectivity offers flexibility but can become dense, opaque, and difficult to control. 
HASTE is designed around this trade-off by separating a sparse physical feasible set from online functional modulation within that set.

\subsection{Efficient Causal Sequence Modeling for Streaming Inference}

Transformers are appealing for long-range temporal modeling, but vanilla self-attention has computation and memory costs that grow with history length \cite{vaswani2017attention}. 
Efficient long-context variants address this bottleneck from complementary angles. 
Transformer-XL introduces segment-level recurrence and cached states to extend effective context \cite{dai2019transformer}, while Longformer restricts attention to local or sparse patterns to improve long-sequence efficiency \cite{Beltagy2020Longformer}. 
These designs are well aligned with streaming deployment, but they are not tailored to EEG fatigue monitoring, where efficient causal temporal reasoning must be coupled with band-aware and topology-constrained spatial modeling. 
Our model builds on this efficient-attention perspective, but integrates it with an explicit physical-graph spatial mechanism and online functional priors for continuous EEG fatigue regression.
\section{Knowledge-Guided Spatial Design and Analysis}
\label{sec:theory}

The preceding discussion suggests that streaming EEG fatigue estimation requires a spatial mechanism that is band-aware, anatomically feasible, stable under online adaptation, and explicitly controllable. We formalize these requirements below.

At each time step $t$, the model receives a multi-band window
$\mathbf{X}_t\in\mathbb{R}^{C\times B}$ and must output a fatigue estimate using only observations up to time $t$.
This setting imposes four constraints on spatial modeling:
the model should remain band-aware as spectral relevance changes over time,
electrode interactions should be anatomically feasible rather than fully dense,
online adaptation should be stable and inspectable to avoid brittle streaming behavior,
and the influence of any spatial prior on attention should be explicit and tunable.

HASTE addresses these requirements with a two-tier design:
a fixed physical graph defines a hard feasible set of spatial interactions,
and a time-varying row-stochastic prior defines a soft, state-dependent preference within that feasible set.
The prior is then injected into neighborhood attention as an explicit logit bias, so its influence is both tunable and transparent.

\subsection{Feasible Neighborhoods and a Graph-Constrained Stochastic Prior}
\label{sec:theory_formulation}

\paragraph{Physical feasibility via a $k$NN electrode graph.}
We encode electrode geometry as a physical graph
$\mathcal{G}^{\mathrm{phy}}=(\mathcal{V},\mathcal{E}^{\mathrm{phy}})$, where $\mathcal{V}$ is the electrode set and
$\mathcal{E}^{\mathrm{phy}}$ connects physically adjacent electrodes (e.g., consistent with standardized montages such as 10--20~\cite{Klem1999TenTwenty}).
For node $i\in\mathcal{V}$, its feasible neighbor set is
$\mathcal{N}(i)=\{j:(i,j)\in\mathcal{E}^{\mathrm{phy}}\}$.
Let $\mathbf{h}_{i,t}$ denote a learnable representation of electrode $i$ before spatial aggregation at time $t$.
We restrict aggregation to the feasible set:
\begin{equation}
\tilde{\mathbf{h}}_{i,t}
=
\sum_{j\in\mathcal{N}(i)}\pi_{ij,t}\,\mathbf{h}_{j,t},
\qquad
\pi_{ij,t}\ge 0,\quad
\sum_{j\in\mathcal{N}(i)}\pi_{ij,t}=1.
\label{eq:neighbor_agg}
\end{equation}
Eq.~\eqref{eq:neighbor_agg} enforces a controllable, sparse spatial receptive field and prevents anatomically implausible mixing.
The remaining question is how to choose $\{\pi_{ij,t}\}$ so that coupling adapts with fatigue state while remaining stable over time.

\paragraph{A row-stochastic coupling prior on the physical graph.}
We encode spatial knowledge at time $t$ as a row-stochastic matrix
$\mathbf{P}_t\in\mathbb{R}^{C\times C}$ supported on $\mathcal{E}^{\mathrm{phy}}$:
\begin{equation}
P_{ij,t}\ge 0,\qquad
\sum_{j=1}^{C}P_{ij,t}=1,\qquad
P_{ij,t}=0\ \text{if } (i,j)\notin\mathcal{E}^{\mathrm{phy}}.
\label{eq:prior_stochastic}
\end{equation}
Each row of $\mathbf{P}_t$ is a probability distribution over $\mathcal{N}(i)$.
Equivalently, $\mathbf{P}_t$ is a Markov transition kernel on $\mathcal{G}^{\mathrm{phy}}$:
$P_{ij,t}$ describes how strongly node $i$ prefers to route information to neighbor $j$ at time $t$.
We do not force features to diffuse according to $\mathbf{P}_t$; instead, we inject it as a bias into attention (Section~\ref{sec:theory_attention}).

\paragraph{Two-tier knowledge: anatomy as a stable baseline, function as adaptive preference.}
We construct $\mathbf{P}_t$ by mixing a static anatomical prior and a dynamic functional prior.
Let $\mathbf{P}^{\mathrm{phy}}$ be a fixed row-stochastic prior supported on $\mathcal{E}^{\mathrm{phy}}$ (geometry-driven baseline).
Let $\mathbf{P}^{\mathrm{func}}_{t,b}$ denote a band-specific functional prior at time $t$ for band $b\in\{1,\ldots,B\}$,
also row-stochastic on $\mathcal{E}^{\mathrm{phy}}$.
Because fatigue effects are frequency-dependent, we fuse band priors using band relevance weights
$\boldsymbol{\alpha}_t=[\alpha_{t,1},\ldots,\alpha_{t,B}]^\top$ with $\alpha_{t,b}\ge 0$ and $\sum_{b=1}^{B}\alpha_{t,b}=1$:
\begin{equation}
\mathbf{P}^{\mathrm{func}}_t
=
\sum_{b=1}^{B}\alpha_{t,b}\mathbf{P}^{\mathrm{func}}_{t,b}.
\label{eq:spectral_fusion}
\end{equation}
We then combine anatomical stability and functional adaptivity with a prior-mixing gate $g_t\in[0,1]$:
\begin{equation}
\mathbf{P}_t
=
g_t\mathbf{P}^{\mathrm{phy}} + (1-g_t)\mathbf{P}^{\mathrm{func}}_t.
\label{eq:gate_fusion}
\end{equation}
Here $g_t$ explicitly controls the balance between anatomy and function:
larger $g_t$ places more weight on $\mathbf{P}^{\mathrm{phy}}$,
whereas smaller $g_t$ allows stronger functional influence.
Since all components are row-stochastic and share the same support, $\mathbf{P}_t$ again satisfies Eq.~\eqref{eq:prior_stochastic}.
Indeed, for any row $i$,
\[
\sum_{j=1}^{C}P_{ij,t}
=
g_t\sum_{j=1}^{C}P^{\mathrm{phy}}_{ij}
+
(1-g_t)\sum_{j=1}^{C}P^{\mathrm{func}}_{ij,t}
=
g_t+(1-g_t)
=
1,
\]
and $P_{ij,t}=0$ whenever $(i,j)\notin\mathcal{E}^{\mathrm{phy}}$ because both
$\mathbf{P}^{\mathrm{phy}}$ and $\mathbf{P}^{\mathrm{func}}_t$ vanish outside the same physical support.
Hence $\mathbf{P}_t$ remains a valid row-stochastic prior on $\mathcal{G}^{\mathrm{phy}}$.
To encourage stable streaming behavior, we later regularize temporal drift of $\mathbf{P}^{\mathrm{func}}_t$ and $g_t$
with $\mathcal{L}_{\mathrm{graph}}$ (Eq.~\eqref{eq:L_graph}).

\subsection{Prior-Biased Attention: Monotonic Control and Drift Stability}
\label{sec:theory_attention}

We convert the prior into an explicit attention bias so that knowledge influence is controllable and inspectable.
For each feasible edge $(i,j)$ with $j\in\mathcal{N}(i)$, let $\tilde{\ell}_{ij,t}$ be a learnable data-driven affinity score.
We define the biased logit
\begin{equation}
\ell_{ij,t}
=
\tilde{\ell}_{ij,t}
+
\beta \log(P_{ij,t}+\epsilon),
\qquad j\in\mathcal{N}(i),
\label{eq:biased_attention}
\end{equation}
where $\beta>0$ controls prior strength and $\epsilon>0$ ensures numerical stability.
Neighborhood softmax yields attention weights
\begin{equation}
\pi_{ij,t}
=
\frac{\exp(\ell_{ij,t})}{\sum_{k\in\mathcal{N}(i)}\exp(\ell_{ik,t})},
\qquad j\in\mathcal{N}(i).
\label{eq:attention_weights}
\end{equation}
Feasibility is guaranteed because attention is computed only over $\mathcal{N}(i)$.
Adaptivity arises because $\mathbf{P}_t$ changes with time and band relevance.
Inspectability follows because $\boldsymbol{\alpha}_t$, $g_t$, and $\mathbf{P}_t$ can be examined directly at inference time.

\begin{proposition}[Posterior form of prior-biased attention]
\label{prop:posterior_form}
Fix node $i$ and time $t$.
Under Eqs.~\eqref{eq:biased_attention}--\eqref{eq:attention_weights},
\begin{equation}
\pi_{ij,t}
=
\frac{\exp(\tilde{\ell}_{ij,t})\,(P_{ij,t}+\epsilon)^{\beta}}
{\sum_{k\in\mathcal{N}(i)}\exp(\tilde{\ell}_{ik,t})\,(P_{ik,t}+\epsilon)^{\beta}},
\qquad j\in\mathcal{N}(i).
\label{eq:posterior_form}
\end{equation}
\end{proposition}

\begin{proof}
For fixed node $i$ and time $t$, substitute Eq.~\eqref{eq:biased_attention} into the numerator of Eq.~\eqref{eq:attention_weights}:
\[
\exp(\ell_{ij,t})
=
\exp\!\bigl(\tilde{\ell}_{ij,t}+\beta\log(P_{ij,t}+\epsilon)\bigr)
=
\exp(\tilde{\ell}_{ij,t})\,\exp\!\bigl(\beta\log(P_{ij,t}+\epsilon)\bigr).
\]
Using $\exp(\beta\log x)=x^\beta$ for $x>0$, we obtain
\[
\exp(\ell_{ij,t})
=
\exp(\tilde{\ell}_{ij,t})(P_{ij,t}+\epsilon)^\beta .
\]
Applying the same expansion to every feasible neighbor $k\in\mathcal{N}(i)$, the normalization term becomes
\[
\sum_{k\in\mathcal{N}(i)}\exp(\ell_{ik,t})
=
\sum_{k\in\mathcal{N}(i)}
\exp(\tilde{\ell}_{ik,t})(P_{ik,t}+\epsilon)^\beta.
\]
Substituting these expressions back into Eq.~\eqref{eq:attention_weights} yields
\[
\pi_{ij,t}
=
\frac{\exp(\tilde{\ell}_{ij,t})(P_{ij,t}+\epsilon)^\beta}
{\sum_{k\in\mathcal{N}(i)}\exp(\tilde{\ell}_{ik,t})(P_{ik,t}+\epsilon)^\beta},
\qquad j\in\mathcal{N}(i),
\]
which is exactly Eq.~\eqref{eq:posterior_form}. Note that $\epsilon>0$ ensures every term is finite and positive, so the neighborhood softmax is well defined even when some prior entries are zero.
\end{proof}

Eq.~\eqref{eq:posterior_form} shows that attention is a normalized product of a data-driven score and a graph prior,
with $\beta$ acting as a direct control knob.
The prior influence is monotonic at the logit level:
\begin{equation}
\frac{\partial \ell_{ij,t}}{\partial P_{ij,t}}
=
\frac{\beta}{P_{ij,t}+\epsilon}
>
0.
\label{eq:monotonic}
\end{equation}
Moreover, holding $\{\tilde{\ell}_{ik,t}\}_{k\in\mathcal{N}(i)}$ and all other prior entries fixed, the induced attention change is also monotonic:
\[
\frac{\partial \pi_{ij,t}}{\partial P_{ij,t}}
=
\pi_{ij,t}(1-\pi_{ij,t})\,\frac{\beta}{P_{ij,t}+\epsilon}
>
0,
\qquad
\frac{\partial \pi_{ik,t}}{\partial P_{ij,t}}
=
-\pi_{ik,t}\pi_{ij,t}\,\frac{\beta}{P_{ij,t}+\epsilon}
<
0
\quad (k\neq j).
\]
Thus increasing the prior weight on one feasible edge increases its own attention share and decreases the competing shares of the remaining neighbors.

\begin{proposition}[Stability under prior drift]
\label{prop:stability}
Fix node $i$ and time $t$, and hold $\{\tilde{\ell}_{ij,t}\}_{j\in\mathcal{N}(i)}$ constant.
Consider two priors $\mathbf{P}_t$ and $\mathbf{P}'_t$ with corresponding neighborhood attention distributions
$\boldsymbol{\pi}_{i,t}$ and $\boldsymbol{\pi}'_{i,t}$ on $\mathcal{N}(i)$.
Let $\mathbf{p}_{i,t}$ and $\mathbf{p}'_{i,t}$ be the $i$-th rows restricted to $\mathcal{N}(i)$.
Then
\begin{equation}
\|\boldsymbol{\pi}_{i,t}-\boldsymbol{\pi}'_{i,t}\|_2
\le
\frac{\beta}{2\epsilon}\,\|\mathbf{p}_{i,t}-\mathbf{p}'_{i,t}\|_2.
\label{eq:stability_bound}
\end{equation}
\end{proposition}

\begin{proof}
Restrict attention to the neighborhood $\mathcal{N}(i)$ and define
\[
\mathbf{u}
=
\tilde{\boldsymbol{\ell}}_{i,t}
+
\mathbf{b}(\mathbf{p}_{i,t}),
\qquad
\mathbf{v}
=
\tilde{\boldsymbol{\ell}}_{i,t}
+
\mathbf{b}(\mathbf{p}'_{i,t}),
\]
where
\[
\mathbf{b}(\mathbf{p})
=
\beta\log(\mathbf{p}+\epsilon\mathbf{1})
\]
is applied elementwise and
$\tilde{\boldsymbol{\ell}}_{i,t}=[\tilde{\ell}_{ij,t}]_{j\in\mathcal{N}(i)}$.
Then
\[
\boldsymbol{\pi}_{i,t}=\mathrm{softmax}(\mathbf{u}),
\qquad
\boldsymbol{\pi}'_{i,t}=\mathrm{softmax}(\mathbf{v}).
\]

We first bound the perturbation of the bias term.
For each coordinate $j\in\mathcal{N}(i)$, the scalar function
$f(x)=\beta\log(x+\epsilon)$ has derivative
\[
f'(x)=\frac{\beta}{x+\epsilon},
\qquad x\ge 0.
\]
Since $x+\epsilon\ge\epsilon$, we have $|f'(x)|\le \beta/\epsilon$ for all $x\ge 0$.
By the mean value theorem,
\[
\bigl|\beta\log(p_{ij,t}+\epsilon)-\beta\log(p'_{ij,t}+\epsilon)\bigr|
\le
\frac{\beta}{\epsilon}\,|p_{ij,t}-p'_{ij,t}|.
\]
Applying this coordinatewise and summing yields
\begin{equation}
\|\mathbf{b}(\mathbf{p}_{i,t})-\mathbf{b}(\mathbf{p}'_{i,t})\|_2
\le
\frac{\beta}{\epsilon}\,
\|\mathbf{p}_{i,t}-\mathbf{p}'_{i,t}\|_2.
\label{eq:bias_lipschitz}
\end{equation}

We next bound the softmax map.
Let $\sigma(\cdot)$ denote softmax over $\mathcal{N}(i)$.
Its Jacobian at a point $\mathbf{z}$ is
\[
\mathbf{J}(\mathbf{z})
=
\mathrm{diag}(\sigma(\mathbf{z}))
-
\sigma(\mathbf{z})\sigma(\mathbf{z})^\top.
\]
Write $\boldsymbol{\rho}=\sigma(\mathbf{z})$.
For row $r$, the absolute row sum is
\[
\sum_{k}|J_{rk}|
=
\rho_r(1-\rho_r)+\sum_{k\neq r}\rho_r\rho_k
=
2\rho_r(1-\rho_r)
\le
\frac{1}{2},
\]
because $\rho_r(1-\rho_r)\le 1/4$.
Hence
\[
\|\mathbf{J}(\mathbf{z})\|_\infty \le \frac{1}{2}.
\]
Since $\mathbf{J}(\mathbf{z})$ is symmetric, $\|\mathbf{J}(\mathbf{z})\|_1=\|\mathbf{J}(\mathbf{z})\|_\infty$, and therefore
\[
\|\mathbf{J}(\mathbf{z})\|_2
\le
\sqrt{\|\mathbf{J}(\mathbf{z})\|_1\|\mathbf{J}(\mathbf{z})\|_\infty}
\le
\frac{1}{2}.
\]
Using the integral form of the mean value theorem,
\[
\sigma(\mathbf{u})-\sigma(\mathbf{v})
=
\int_0^1
\mathbf{J}\!\bigl(\mathbf{v}+\tau(\mathbf{u}-\mathbf{v})\bigr)
(\mathbf{u}-\mathbf{v})\,d\tau .
\]
Taking $\ell_2$ norms and using the uniform Jacobian bound gives
\begin{equation}
\|\mathrm{softmax}(\mathbf{u})-\mathrm{softmax}(\mathbf{v})\|_2
\le
\frac{1}{2}\|\mathbf{u}-\mathbf{v}\|_2 .
\label{eq:softmax_lipschitz}
\end{equation}

Finally, $\tilde{\boldsymbol{\ell}}_{i,t}$ is held fixed, so
\[
\mathbf{u}-\mathbf{v}
=
\mathbf{b}(\mathbf{p}_{i,t})-\mathbf{b}(\mathbf{p}'_{i,t}).
\]
Combining Eqs.~\eqref{eq:bias_lipschitz} and \eqref{eq:softmax_lipschitz}, we obtain
\[
\|\boldsymbol{\pi}_{i,t}-\boldsymbol{\pi}'_{i,t}\|_2
\le
\frac{1}{2}
\|\mathbf{b}(\mathbf{p}_{i,t})-\mathbf{b}(\mathbf{p}'_{i,t})\|_2
\le
\frac{\beta}{2\epsilon}\,
\|\mathbf{p}_{i,t}-\mathbf{p}'_{i,t}\|_2,
\]
which is Eq.~\eqref{eq:stability_bound}.
\end{proof}

Proposition~\ref{prop:stability} turns the streaming stability requirement into a concrete design target:
controlling the drift of $\mathbf{P}_t$ directly controls the prior-induced perturbation of neighborhood attention.
This motivates regularizing temporal drift of $\mathbf{P}^{\mathrm{func}}_t$ and $g_t$ during training (Eq.~\eqref{eq:L_graph}).

\begin{lemma}[Gate-induced drift decomposition]
\label{lem:gate_drift}
For $\mathbf{P}_t$ in Eq.~\eqref{eq:gate_fusion},
\begin{equation}
\mathbf{P}_t-\mathbf{P}_{t-1}
=
(g_t-g_{t-1})(\mathbf{P}^{\mathrm{phy}}-\mathbf{P}^{\mathrm{func}}_{t-1})
+
(1-g_t)(\mathbf{P}^{\mathrm{func}}_{t}-\mathbf{P}^{\mathrm{func}}_{t-1}).
\label{eq:gate_drift}
\end{equation}
\end{lemma}

\begin{proof}
Starting from Eq.~\eqref{eq:gate_fusion},
\[
\mathbf{P}_t
=
g_t\mathbf{P}^{\mathrm{phy}}+(1-g_t)\mathbf{P}^{\mathrm{func}}_t,
\qquad
\mathbf{P}_{t-1}
=
g_{t-1}\mathbf{P}^{\mathrm{phy}}+(1-g_{t-1})\mathbf{P}^{\mathrm{func}}_{t-1}.
\]
Subtracting the two expressions gives
\[
\mathbf{P}_t-\mathbf{P}_{t-1}
=
(g_t-g_{t-1})\mathbf{P}^{\mathrm{phy}}
+
(1-g_t)\mathbf{P}^{\mathrm{func}}_t
-
(1-g_{t-1})\mathbf{P}^{\mathrm{func}}_{t-1}.
\]
To separate gate variation from functional-prior variation, add and subtract
$(1-g_t)\mathbf{P}^{\mathrm{func}}_{t-1}$:
\[
\mathbf{P}_t-\mathbf{P}_{t-1}
=
(g_t-g_{t-1})\mathbf{P}^{\mathrm{phy}}
+
(1-g_t)\mathbf{P}^{\mathrm{func}}_t
-
(1-g_t)\mathbf{P}^{\mathrm{func}}_{t-1}
+
(1-g_t)\mathbf{P}^{\mathrm{func}}_{t-1}
-
(1-g_{t-1})\mathbf{P}^{\mathrm{func}}_{t-1}.
\]
Grouping terms yields
\[
\mathbf{P}_t-\mathbf{P}_{t-1}
=
(g_t-g_{t-1})\bigl(\mathbf{P}^{\mathrm{phy}}-\mathbf{P}^{\mathrm{func}}_{t-1}\bigr)
+
(1-g_t)\bigl(\mathbf{P}^{\mathrm{func}}_{t}-\mathbf{P}^{\mathrm{func}}_{t-1}\bigr),
\]
which is Eq.~\eqref{eq:gate_drift}.
\end{proof}

Eq.~\eqref{eq:gate_drift} separates two sources of change:
the first term measures how the model shifts its balance between anatomical and functional priors (gate drift),
and the second term measures how quickly the functional prior itself evolves,
weighted by the current functional contribution $(1-g_t)$.

\begin{corollary}[Gate-aware bound on prior-induced attention perturbation]
\label{cor:attention_drift}
Fix a node $i$ and a data-driven affinity vector
$\tilde{\boldsymbol{\ell}}_{i}=[\tilde{\ell}_{ij}]_{j\in\mathcal{N}(i)}$.
Let $\boldsymbol{\pi}_{i}(\mathbf{P})$ denote the neighborhood attention distribution produced by
Eqs.~\eqref{eq:biased_attention}--\eqref{eq:attention_weights} when the prior row is taken from $\mathbf{P}$ and $\tilde{\boldsymbol{\ell}}_{i}$ is held fixed.
For two consecutive priors $\mathbf{P}_{t-1}$ and $\mathbf{P}_{t}$ defined by Eq.~\eqref{eq:gate_fusion},
\begin{equation}
\|\boldsymbol{\pi}_{i}(\mathbf{P}_{t})-\boldsymbol{\pi}_{i}(\mathbf{P}_{t-1})\|_2
\le
\frac{\beta}{2\epsilon}
\left[
|g_t-g_{t-1}|\,\|\mathbf{p}^{\mathrm{phy}}_{i}-\mathbf{p}^{\mathrm{func}}_{i,t-1}\|_2
+
(1-g_t)\,\|\mathbf{p}^{\mathrm{func}}_{i,t}-\mathbf{p}^{\mathrm{func}}_{i,t-1}\|_2
\right],
\label{eq:attention_drift_decomp}
\end{equation}
where
$\mathbf{p}^{\mathrm{phy}}_{i}=\mathbf{P}^{\mathrm{phy}}(i,\mathcal{N}(i))$
and
$\mathbf{p}^{\mathrm{func}}_{i,t}=\mathbf{P}^{\mathrm{func}}_{t}(i,\mathcal{N}(i))$.
\end{corollary}

\begin{proof}
Restrict Eq.~\eqref{eq:gate_drift} to row $i$ and to the neighborhood coordinates $\mathcal{N}(i)$.
This gives
\[
\mathbf{p}_{i,t}-\mathbf{p}_{i,t-1}
=
(g_t-g_{t-1})(\mathbf{p}^{\mathrm{phy}}_{i}-\mathbf{p}^{\mathrm{func}}_{i,t-1})
+
(1-g_t)(\mathbf{p}^{\mathrm{func}}_{i,t}-\mathbf{p}^{\mathrm{func}}_{i,t-1}),
\]
where $\mathbf{p}_{i,t}=\mathbf{P}_t(i,\mathcal{N}(i))$.
Applying the triangle inequality yields
\[
\|\mathbf{p}_{i,t}-\mathbf{p}_{i,t-1}\|_2
\le
|g_t-g_{t-1}|\,\|\mathbf{p}^{\mathrm{phy}}_{i}-\mathbf{p}^{\mathrm{func}}_{i,t-1}\|_2
+
(1-g_t)\,\|\mathbf{p}^{\mathrm{func}}_{i,t}-\mathbf{p}^{\mathrm{func}}_{i,t-1}\|_2.
\]
Now apply Proposition~\ref{prop:stability} with the data-driven affinity vector fixed at $\tilde{\boldsymbol{\ell}}_{i}$ to obtain
\[
\|\boldsymbol{\pi}_{i}(\mathbf{P}_{t})-\boldsymbol{\pi}_{i}(\mathbf{P}_{t-1})\|_2
\le
\frac{\beta}{2\epsilon}\,
\|\mathbf{p}_{i,t}-\mathbf{p}_{i,t-1}\|_2.
\]
Combining the two inequalities proves Eq.~\eqref{eq:attention_drift_decomp}.
\end{proof}

Corollary~\ref{cor:attention_drift} makes explicit how the two regularized quantities in $\mathcal{L}_{\mathrm{graph}}$ affect the prior-induced perturbation of attention.
Suppressing gate drift controls how rapidly the model shifts between anatomical and functional preferences,
whereas suppressing $\mathbf{P}^{\mathrm{func}}_t-\mathbf{P}^{\mathrm{func}}_{t-1}$ controls how quickly the functional prior itself changes.
This is precisely why regularizing both terms promotes stable streaming behavior while preserving adaptive spatial coupling.
\section{Method: HASTE Architecture}
\label{sec:method}

HASTE instantiates the spatial design above in a causal spatio-temporal architecture for subject-independent continuous fatigue regression. At time step $t$, the model receives one multi-band window $\mathbf{X}_t\in\mathbb{R}^{C\times B}$ and predicts $\hat{y}_t\in[0,1]$ using only observations up to time $t$. Each update has four stages: online construction of a band-aware spatial prior on the physical graph, prior-biased spatial encoding of the current window, bounded causal temporal aggregation, and regression of the current fatigue state. Figure~2 summarizes this pipeline. The upper branch constructs the fused spatial prior through causal spectral fusion, band-wise functional priors, active-band mixing, and anatomical--functional fusion, whereas the lower branch uses the resulting prior for topology-constrained spatial attention, tokenization, causal local temporal attention, and prediction.

\begin{figure*}[htbp]
\centering
\includegraphics[width=0.98\linewidth]{Figure2.pdf}
\caption{Detailed framework of HASTE. The upper branch builds the knowledge prior from the current multi-band window via causal spectral fusion, a physical topology prior, band-wise functional priors, Top-$K$ active-band mixing, and final gate-based prior fusion. The lower branch uses the resulting prior to guide masked spatial attention, tokenization, local temporal attention, and final fatigue regression.}
\end{figure*}


\subsection{Streaming Setup and Persistent States}
\label{sec:method_setup}

We use $C$ electrodes and $B$ frequency bands for each of $T$ windows.
Let $\mathbf{X}\in\mathbb{R}^{C\times T\times B}$ be the full input tensor, where $X_{c,t,b}$ is the feature at electrode $c$,
time $t$, and band $b$.
At time $t$, the model observes $\mathbf{X}_t\triangleq\mathbf{X}_{:,t,:}\in\mathbb{R}^{C\times B}$.
Let $y_t\in[0,1]$ be the normalized fatigue label and $\hat{y}_t\in[0,1]$ the prediction.
The model is causal:
\begin{equation}
\hat{y}_t = f(\mathbf{X}_{1:t};\theta).
\label{eq:causal_prediction}
\end{equation}

Streaming inference relies on the continuous maintenance of three persistent states. Specifically, the process tracks a spectral state $\mathbf{r}_{t-1}\in\mathbb{R}^{d_r}$ alongside a temporal key-value cache $\mathcal{C}_{t-1}$, which stores the most recent $L-1$ key and value pairs per attention head. Furthermore, the system preserves a static physical graph $\mathcal{G}^{\mathrm{phy}}=(\mathcal{V},\mathcal{E}^{\mathrm{phy}})$ in conjunction with an anatomical prior $\mathbf{P}^{\mathrm{phy}}$.

\subsection{Build the Knowledge Prior Online}
\label{sec:method_prior}

This block instantiates the band-aware stochastic prior in Section~\ref{sec:theory_formulation} under streaming and compute constraints.
\paragraph{Causal spectral fusion.}
We estimate time-varying band relevance weights $\boldsymbol{\alpha}_t\in\mathbb{R}^{B}$.
We compute a global band descriptor $\mathbf{u}_t\in\mathbb{R}^{B}$ by averaging $\mathbf{X}_t$ over electrodes:
$u_{t,b}=\frac{1}{C}\sum_{c=1}^{C}X_{c,t,b}$.
A causal GRU updates the spectral state $\mathbf{r}_t=\mathrm{GRU}(\mathbf{u}_t,\mathbf{r}_{t-1})$~\cite{Cho2014GRU},
and we compute
$\boldsymbol{\alpha}_t=\mathrm{Softmax}(\mathbf{W}_s\mathbf{r}_t+\mathbf{b}_s)$.
We also fuse multi-band features into one scalar per electrode:
$x_{c,t}=\sum_{b=1}^{B}\alpha_{t,b}X_{c,t,b}$,
and collect $\mathbf{x}_t=[x_{1,t},\ldots,x_{C,t}]^\top\in\mathbb{R}^{C}$.
To stabilize streaming behavior, we regularize spectral drift:
\begin{equation}
\mathcal{L}_{\mathrm{band}}
=
\frac{1}{T-1}\sum_{t=2}^{T}\|\boldsymbol{\alpha}_{t}-\boldsymbol{\alpha}_{t-1}\|_2^2.
\label{eq:L_band}
\end{equation}

\paragraph{Band gating for compute-bounded functional priors.}
Constructing $\mathbf{P}^{\mathrm{func}}_{t,b}$ for all $B$ bands can dominate spatial cost.
We therefore select an active band set $\mathcal{B}_t$ of size $K$ via $\mathcal{B}_t=\mathrm{TopK}(\boldsymbol{\alpha}_t,K)$.
Let $m_{t,b}\in\{0,1\}$ indicate whether band $b$ is active.
We form masked-and-renormalized weights
\begin{equation}
\bar{\alpha}_{t,b}
=
\frac{m_{t,b}\alpha_{t,b}}{\sum_{b'\in\mathcal{B}_t}\alpha_{t,b'}},
\qquad
\sum_{b\in\mathcal{B}_t}\bar{\alpha}_{t,b}=1.
\label{eq:masked_weights}
\end{equation}
Optionally, during training we replace the hard mask with a differentiable approximation
$\tilde{\mathbf{m}}_t=\mathrm{SOFTTopK}(\boldsymbol{\alpha}_t;K,\tau)\in[0,1]^B$~\cite{xie2020differentiable}
to reduce train--test mismatch.

\paragraph{Physical graph and anatomical prior.}
We instantiate the physical graph using a $k$NN rule on electrode coordinates, symmetrize edges, and add self-loops.
Let $\mathbf{A}^{\mathrm{phy}}\in\{0,1\}^{C\times C}$ be the adjacency.
The anatomical prior is row-normalized adjacency:
\begin{equation}
P^{\mathrm{phy}}_{ij}
=
\frac{A^{\mathrm{phy}}_{ij}}{\sum_{k=1}^{C}A^{\mathrm{phy}}_{ik}}.
\label{eq:phy_norm}
\end{equation}

\paragraph{Band-wise functional prior (only on active bands).}
\label{sec:method_func}
For each electrode $c$ and band $b$, we embed the scalar feature into $\mathbb{R}^{d}$:
\begin{equation}
\mathbf{e}_{c,t,b}
=
\mathbf{W}^{(e)}_b X_{c,t,b}+\mathbf{b}^{(e)}_b
\in \mathbb{R}^{d}.
\label{eq:band_embed}
\end{equation}
To reduce compute, we measure similarity in a lower dimension $d_s\ll d$.
For each feasible edge $(i,j)$ with $j\in\mathcal{N}(i)$:
\begin{equation}
s_{ij,t,b}
=
\frac{\big(\mathbf{W}_q\mathbf{e}_{i,t,b}\big)^\top\big(\mathbf{W}_k\mathbf{e}_{j,t,b}\big)}{\sqrt{d_s}},
\qquad
\mathbf{W}_q,\mathbf{W}_k\in\mathbb{R}^{d_s\times d}.
\label{eq:similarity_small}
\end{equation}
We obtain a row-stochastic prior by neighborhood softmax:
\begin{equation}
P^{\mathrm{func}}_{ij,t,b}
=
\begin{cases}
\displaystyle\frac{\exp(s_{ij,t,b})}{\sum_{k\in\mathcal{N}(i)}\exp(s_{ik,t,b})}, & j\in\mathcal{N}(i),\\[6pt]
0, & \text{otherwise}.
\end{cases}
\label{eq:P_func_band}
\end{equation}
In streaming inference, we only construct $\mathbf{P}^{\mathrm{func}}_{t,b}$ for $b\in\mathcal{B}_t$.
\paragraph{Band mixing and prior-mixing gate.}
We fuse active-band functional priors with masked weights:
\begin{equation}
\mathbf{P}^{\mathrm{func}}_t
=
\sum_{b\in\mathcal{B}_t}\bar{\alpha}_{t,b}\mathbf{P}^{\mathrm{func}}_{t,b}.
\label{eq:masked_fusion}
\end{equation}
To remain consistent with Top-$K$ computation, we predict the prior-mixing gate using active-band embeddings only:
\begin{equation}
\bar{\mathbf{e}}_t
=
\frac{1}{C|\mathcal{B}_t|}\sum_{c=1}^{C}\sum_{b\in\mathcal{B}_t}\mathbf{e}_{c,t,b},
\qquad
g_t
=
\sigma\!\left(\mathbf{w}_g^\top\bar{\mathbf{e}}_t+b_g\right).
\label{eq:gate}
\end{equation}
The final knowledge prior is the convex mix
\begin{equation}
\mathbf{P}_t
=
g_t\mathbf{P}^{\mathrm{phy}} + (1-g_t)\mathbf{P}^{\mathrm{func}}_t.
\label{eq:gate_fusion_method}
\end{equation}
where larger $g_t$ corresponds to a stronger anatomical contribution,
and smaller $g_t$ corresponds to a stronger functional contribution.

\subsection{Encode and Predict with Bounded Streaming State}
\label{sec:method_encode}

\paragraph{Prior-biased spatial attention on the physical graph.}
We perform graph-restricted multi-head attention (GAT-style)~\cite{velivckovic2017graph}.
We lift the fused scalar to a node representation
\begin{equation}
\mathbf{h}^{(0)}_{i,t}
=
\mathbf{w}_{\mathrm{in}} x_{i,t} + \mathbf{b}_{\mathrm{in}}
\in \mathbb{R}^{d}.
\label{eq:node_embed}
\end{equation}
We parameterize prior strength as $\beta=\log(1+\exp(\tilde{\beta}))$ so $\beta>0$.
For spatial layer $l$ and head $h$, the logit is
\begin{equation}
\ell^{(h)}_{ij,t}
=
\mathrm{LeakyReLU}\!\left(
\mathbf{a}^{(h)\top}
\left[
\mathbf{W}^{(h)}\mathbf{h}^{(l)}_{i,t}\,\|\,\mathbf{W}^{(h)}\mathbf{h}^{(l)}_{j,t}
\right]
\right)
+
\beta\log(P_{ij,t}+\epsilon),
\qquad j\in\mathcal{N}(i),
\label{eq:spatial_logit}
\end{equation}
with $\epsilon=10^{-6}$.
Neighborhood softmax yields $\pi^{(h)}_{ij,t}$ and aggregation is
\begin{equation}
\tilde{\mathbf{h}}^{(l,h)}_{i,t}
=
\sum_{j\in\mathcal{N}(i)}\pi^{(h)}_{ij,t}\mathbf{W}^{(h)}\mathbf{h}^{(l)}_{j,t}.
\label{eq:spatial_agg}
\end{equation}
Heads are concatenated and activated:
\begin{equation}
\mathbf{h}^{(l+1)}_{i,t}
=
\mathrm{ELU}\!\left(\mathrm{Concat}_{h=1}^{H_s}\tilde{\mathbf{h}}^{(l,h)}_{i,t}\right).
\label{eq:head_concat}
\end{equation}
After $L_s$ layers, we obtain $\mathbf{H}_{\mathrm{spatial},t}\in\mathbb{R}^{C\times d}$.

\paragraph{Tokenization and causal local temporal attention.}
We flatten $\mathbf{H}_{\mathrm{spatial},t}$ into a temporal token
\begin{equation}
\mathbf{z}_t
=
\mathbf{W}_{\mathrm{proj}}\,\mathrm{vec}(\mathbf{H}_{\mathrm{spatial},t})
+
\mathbf{b}_{\mathrm{proj}}
\in\mathbb{R}^{d_{\mathrm{model}}}.
\label{eq:temporal_token}
\end{equation}
We apply \emph{causal local temporal attention} (CLTA)~\cite{vaswani2017attention}, restricting attention to the most recent $L$ windows with mask
\begin{equation}
m_{t,s}
=
\begin{cases}
0, & \max(1,t-L+1)\le s\le t,\\
-\infty, & \text{otherwise}.
\end{cases}
\label{eq:causal_mask}
\end{equation}
In streaming mode, we cache only the most recent $L-1$ key/value pairs per head, so the temporal state does not grow with $T$.
(We cache both keys and values; constants are omitted in big-$O$ notation.)

\paragraph{Prediction head.}
Let $\mathbf{h}_t\in\mathbb{R}^{d_{\mathrm{model}}}$ be the temporal output at time $t$.
We predict fatigue with a sigmoid output:
\begin{equation}
\hat{y}_t
=
\sigma(\mathbf{w}_{\mathrm{out}}^{\top}\mathbf{h}_t + b_{\mathrm{out}}).
\label{eq:prediction}
\end{equation}

\paragraph{Streaming forward pass.}
Algorithm~\ref{alg:haste_stream} summarizes streaming inference at time $t$.

\begin{algorithm}[!t]
\caption{Streaming inference of HASTE at time $t$}
\label{alg:haste_stream}
\begin{algorithmic}
\Require Window input $\mathbf{X}_t\in\mathbb{R}^{C\times B}$; physical graph $\mathcal{G}^{\mathrm{phy}}$ and $\mathbf{P}^{\mathrm{phy}}$;
previous spectral state $\mathbf{r}_{t-1}$; temporal cache $\mathcal{C}_{t-1}$; hyperparameters $K$ and $L$
\Ensure Prediction $\hat{y}_t$; updated spectral state $\mathbf{r}_{t}$; updated cache $\mathcal{C}_{t}$
\State Compute $\mathbf{u}_t$ by averaging $\mathbf{X}_t$ over electrodes
\State Update $\mathbf{r}_t\leftarrow \mathrm{GRU}(\mathbf{u}_t,\mathbf{r}_{t-1})$ and compute $\boldsymbol{\alpha}_t$
\State Select $\mathcal{B}_t\leftarrow \mathrm{TopK}(\boldsymbol{\alpha}_t,K)$ and masked weights $\{\bar{\alpha}_{t,b}\}_{b\in\mathcal{B}_t}$
\For{each $b\in\mathcal{B}_t$}
\State Compute embeddings $\mathbf{e}_{:,t,b}$ by Eq.~\eqref{eq:band_embed}
\State Compute similarities on physical edges by Eq.~\eqref{eq:similarity_small}
\State Compute $\mathbf{P}^{\mathrm{func}}_{t,b}$ by Eq.~\eqref{eq:P_func_band}
\EndFor
\State Fuse $\mathbf{P}^{\mathrm{func}}_t$ by Eq.~\eqref{eq:masked_fusion}
\State Compute $g_t$ by Eq.~\eqref{eq:gate} and obtain $\mathbf{P}_t$ by Eq.~\eqref{eq:gate_fusion_method}
\State Fuse features to $\mathbf{x}_t$ using $\boldsymbol{\alpha}_t$
\State Run spatial encoder to obtain $\mathbf{H}_{\mathrm{spatial},t}$
\State Tokenize to $\mathbf{z}_t$ by Eq.~\eqref{eq:temporal_token}
\State Apply CLTA with mask in Eq.~\eqref{eq:causal_mask} and update cache $\mathcal{C}_t$
\State Predict $\hat{y}_t$ by Eq.~\eqref{eq:prediction}
\end{algorithmic}
\end{algorithm}

\subsection{Objective and Streaming Cost}
\label{sec:method_objective}


\paragraph{Training objective.}
We train with regression accuracy, output smoothness, and drift regularization:
\begin{equation}
\mathcal{L}
=
\mathcal{L}_{\mathrm{reg}}
+
\lambda_{\mathrm{pred}}\mathcal{L}_{\mathrm{pred}}
+
\lambda_{\mathrm{band}}\mathcal{L}_{\mathrm{band}}
+
\lambda_{\mathrm{graph}}\mathcal{L}_{\mathrm{graph}}.
\label{eq:final_loss}
\end{equation}
The regression loss is
\begin{equation}
\mathcal{L}_{\mathrm{reg}}
=
\frac{1}{T}\sum_{t=1}^{T}(y_t-\hat{y}_t)^2,
\label{eq:L_reg}
\end{equation}
and the prediction smoothness loss is
\begin{equation}
\mathcal{L}_{\mathrm{pred}}
=
\frac{1}{T-1}\sum_{t=2}^{T}(\hat{y}_{t}-\hat{y}_{t-1})^2.
\label{eq:L_pred}
\end{equation}
We penalize drift of the functional prior and the prior-mixing gate:
\begin{equation}
\mathcal{L}_{\mathrm{graph}}
=
\frac{1}{T-1}\sum_{t=2}^{T}\left\|\mathbf{P}^{\mathrm{func}}_{t}-\mathbf{P}^{\mathrm{func}}_{t-1}\right\|_{F}^{2}
+
\frac{1}{T-1}\sum_{t=2}^{T}(g_{t}-g_{t-1})^2.
\label{eq:L_graph}
\end{equation}

\paragraph{Streaming cost.}
Let $\bar{d}$ be the average node degree in the physical graph, so $|\mathcal{E}^{\mathrm{phy}}|=O(C\bar{d})$.
With band gating, functional prior construction is performed for $K$ bands and only on physical edges.
Using the low-dimensional similarity in Eq.~\eqref{eq:similarity_small}, its per-window cost is $O(KC\bar{d}d_s)$.
Graph-restricted spatial attention costs $O(H_sC\bar{d}d)$.
CLTA costs $O(HLd_k)$ per window.

We summarize per-window complexity as
\begin{equation}
\mathcal{C}_{\mathrm{win}}
=
O\!\left(KC\bar{d}d_s + H_sC\bar{d}d + HLd_k\right),
\label{eq:complexity_win}
\end{equation}
and persistent memory as
\begin{equation}
\mathcal{M}_{\mathrm{state}}
=
O\!\left(d_r + HLd_k + |\mathcal{E}^{\mathrm{phy}}|\right),
\label{eq:memory_state}
\end{equation}
which does not grow with session length $T$.
\section{Experiments}
\label{sec:experiments}


Our experiments assess whether HASTE can serve as a deployable streaming fatigue estimator.
We organize the evaluation around three questions that form a chain of evidence:
(1) \emph{Accuracy:} can HASTE achieve accurate \emph{subject-independent} continuous regression under strict causality?
(2) \emph{Deployability:} does it maintain bounded latency and memory as session length grows?
(3) \emph{Mechanisms:} which components are responsible for the gains, and what internal signals explain the behavior?
We answer these questions using a unified protocol that controls inputs, splits, and causality across all methods.

\subsection{Evaluation Setup}
\label{sec:exp_setup}

\paragraph{Datasets.}
We evaluate HASTE on three public EEG fatigue datasets that differ in supervision source, montage density, and release format (Table~\ref{tab:dataset_summary}). Together, they span heterogeneous targets---continuous PERCLOS trajectories, reaction-time-derived indices, and physician-annotated fatigue levels mapped to $[0,1]$---and cover increasingly dense montages, which stress-test topology-guided spatial modeling across settings.

\emph{SEED-VIG}~\cite{Zheng2017CriticalBands} provides a public feature-level release for 23 subjects performing a prolonged simulated driving task, together with continuous PERCLOS labels in $[0,1]$. In this work, we use its released five-band DE features rather than claiming a unified raw-signal preprocessing pipeline for this dataset.

\emph{SADT}~\cite{cao2019multi} contains 62 sessions from 27 subjects recorded during a 90-minute sustained-attention driving task. The public raw release contains 32 recorded signals sampled at 500\,Hz, including 30 scalp EEG channels and two mastoid references; in our implementation, we retain the 30 scalp EEG channels. Lane-departure events are induced and the corresponding reaction time (RT) is provided. We convert RT into a continuous drowsiness index (DI) in $[0,1]$ via a monotonic, long-tail-robust mapping.

\emph{MPD-DF}~\cite{Li2026MPDDF} provides 32-channel EEG from 50 subjects in a standardized 2-hour simulated driving task, together with physician annotations of fatigue levels based on EEG and multimodal physiological evidence. The public release provides second-resolved aligned labels derived from five ordered fatigue levels; in our regression setting, we map these ordered labels to $[0,1]$.

\paragraph{Streaming protocol and pre-processing.}
All methods are evaluated under strict causality: at each time step $t$, predictions may use only observations up to and including $t$. For attention-based baselines with global self-attention (e.g., EEG-Conformer), we enforce non-anticipativity using a standard causal mask so each query attends only to past tokens. This preserves the original full-history modeling form, but the per-update compute/state can still grow with the session length $T$. In contrast, HASTE uses causal local temporal attention with a fixed look-back window and a key--value cache, so its persistent state does not grow with $T$.

For consistency across datasets, we operate on 1-second updates represented by five-band DE features. For SADT and MPD-DF, raw EEG is band-pass filtered to 1--50\,Hz, downsampled to 128\,Hz, segmented into non-overlapping 1-second windows, and converted to DE features in the five commonly used bands $\delta$ (1--3\,Hz), $\theta$ (4--7\,Hz), $\alpha$ (8--13\,Hz), $\beta$ (14--30\,Hz), and $\gamma$ (31--50\,Hz). For SEED-VIG, we use the public five-band DE feature release and align it to the same streaming evaluation protocol. Continuous or mapped labels are assigned to each 1-second update, yielding $\mathbf{X}_t \in \mathbb{R}^{C \times 5}$ at every time step.

Throughout, streaming refers to step-wise model updates on sequential window features: the model does not access future windows during inference. Unless stated otherwise, the reported per-update latency and GPU memory profile the neural model forward pass on precomputed features (batch size 1) and exclude feature extraction.

\begin{table}[!t]
\centering
\caption{Summary of the three evaluation datasets.}
\label{tab:dataset_summary}
\renewcommand{\arraystretch}{1.12}
\setlength{\tabcolsep}{3pt} 
\resizebox{\columnwidth}{!}{
\begin{tabular}{lccc}
\toprule
& \textbf{SEED-VIG} & \textbf{SADT} & \textbf{MPD-DF} \\
\midrule
Publication year          & 2017       & 2019       & 2026 \\
Subjects / sessions       & 23\,/\,23  & 27\,/\,62  & 50\,/\,50 \\
Used channels $C$         & 17         & 30         & 32 \\
Native release            & DE features + PERCLOS & Raw EEG + RT & Raw EEG + physician labels \\
Raw EEG sampling rate (Hz)& --         & 500        & 500 \\
Session duration          & ${\sim}$2\,h & ${\sim}$90\,min & ${\sim}$2\,h \\
Target source             & PERCLOS    & DI (from RT) & Ordered physician levels \\
Regression target used    & $[0,1]$    & $[0,1]$    & mapped to $[0,1]$ \\
Evaluation protocol       & LOSO       & LOSO       & LOSO \\
\bottomrule
\end{tabular}%
}
\end{table}
\paragraph{Baselines.}
We compare HASTE with six representative baselines spanning CNNs, RNNs, Transformers, and GNNs.
All baselines are re-implemented under the same streaming protocol, the same DE features, and the same LOSO splits.
\begin{itemize}
  \item \textbf{EEGNet}~\cite{Lawhern2018EEGNet}: compact temporal--spatial CNN adapted to regression with causal padding.
  \item \textbf{TSception}~\cite{ding2023tsception}: multi-scale temporal--spatial CNN adapted to regression with causal convolutions.
  \item \textbf{EEG-Conformer}~\cite{song2023eeg}: CNN--Transformer hybrid with full-history causal attention (causal mask).
  \item \textbf{RGNN}~\cite{zhong2020rgnn}: regularized GNN with a biologically informed adjacency; we use a regression head.
  \item \textbf{NHGNet}~\cite{xu2025nhgnet}: GNN with a learnable dense graph module; we use a regression head.
  \item \textbf{Compact GRU}: two-layer GRU with a linear regression head, serving as a minimal bounded-memory streaming baseline.
\end{itemize}

\paragraph{Implementation details.}
Table~\ref{tab:hyperparams} summarizes shared hyper-parameters.
We set the physical $k$NN degree to $k{=}5$ for SEED-VIG (17 channels) and $k{=}7$ for SADT/MPD-DF (30/32 channels),
reflecting higher montage density.
Unless stated otherwise, experiments run on a single NVIDIA RTX 3090 GPU.

\begin{table}[!t]
\centering
\caption{Implementation details and hyper-parameters (shared unless otherwise stated).}
\label{tab:hyperparams}
\renewcommand{\arraystretch}{1.10}
\setlength{\tabcolsep}{6pt}
\begin{tabular}{ll}
\toprule
\textbf{Item} & \textbf{Value} \\
\midrule
Pre-processing & Band-pass 1--50\,Hz; downsample to 128\,Hz \\
Windowing & 1\,s non-overlapping windows; labels interpolated to centers \\
Features & DE in 5 bands ($\delta,\theta,\alpha,\beta,\gamma$) \\
Topology degree $k$ & 5 (SEED-VIG); 7 (SADT/MPD-DF) \\
Look-back $L$ (CLTA) & 10 \\
Active bands $K$ & 3 \\
Similarity dim $d_s$ & 16 \\
Dims & $d{=}64$, $d_{\mathrm{model}}{=}128$, $d_r{=}32$ \\
Spatial module & $H_s{=}4$ heads, $L_s{=}2$ layers \\
Temporal module & $H{=}4$ heads \\
Optimizer & AdamW ($lr{=}1\mathrm{e}{-}3$, $wd{=}1\mathrm{e}{-}4$) \\
Training & batch size 32; max epochs 100; 3 seeds \\
Loss weights & $\lambda_{\mathrm{pred}}{=}0.1$, $\lambda_{\mathrm{band}}{=}0.01$, $\lambda_{\mathrm{graph}}{=}0.01$ \\
\bottomrule
\end{tabular}
\end{table}

\subsection{Main Results}
\label{sec:exp_main}

Table~\ref{tab:main_results} reports subject-independent continuous fatigue regression under streaming inference on all three datasets.
HASTE achieves the best PCC and the lowest RMSE/MAE on each dataset.
The gains are consistent across heterogeneous label sources and montage densities, supporting robust cross-subject generalization
in realistic streaming settings.


\begin{table}[!t]
\centering
\caption{Subject-independent continuous fatigue regression under streaming inference (LOSO, mean\,$\pm$\,std over 3 seeds). Best and second best per dataset/metric are in \textbf{bold} and \underline{underlined}. Cost denotes per-update state/memory scaling with respect to session length $T$.}
\label{tab:main_results}

\renewcommand{\arraystretch}{1.15}
\setlength{\tabcolsep}{4.2pt}
\footnotesize

\begin{tabular}{@{}l c c r r r@{}}
\toprule
\multirow{2}{*}{\textbf{Method}} &
\multirow{2}{*}{\textbf{Paradigm}} &
\multirow{2}{*}{\textbf{Cost}} &
\multicolumn{3}{c}{\textbf{Metrics}} \\
\cmidrule(lr){4-6}
& & & \textbf{PCC}$\uparrow$ & \textbf{RMSE}$\downarrow$ & \textbf{MAE}$\downarrow$ \\
\midrule

\multicolumn{6}{@{}l@{}}{\textit{SEED-VIG}} \\
EEGNet        & CNN           & $O(1)$ & 0.6847\,${\pm}$\,0.018 & 0.1583\,${\pm}$\,0.007 & 0.1221\,${\pm}$\,0.006 \\
TSception     & CNN           & $O(1)$ & 0.7159\,${\pm}$\,0.015 & 0.1479\,${\pm}$\,0.006 & 0.1138\,${\pm}$\,0.005 \\
EEG-Conformer & CNN+Trf       & $O(T)$ & \underline{0.7831\,${\pm}$\,0.012} & \underline{0.1322\,${\pm}$\,0.005} & \underline{0.1015\,${\pm}$\,0.004} \\
RGNN          & GNN (fixed)   & $O(1)$ & 0.7392\,${\pm}$\,0.016 & 0.1418\,${\pm}$\,0.006 & 0.1093\,${\pm}$\,0.005 \\
NHGNet        & GNN (learned) & $O(1)$ & 0.7634\,${\pm}$\,0.014 & 0.1371\,${\pm}$\,0.006 & 0.1058\,${\pm}$\,0.005 \\
Compact GRU   & RNN           & $O(1)$ & 0.6518\,${\pm}$\,0.021 & 0.1647\,${\pm}$\,0.008 & 0.1274\,${\pm}$\,0.007 \\
\addlinespace[2pt]
\rowcolor{blue!7}
\textbf{HASTE (ours)} & Topo+Attn & $O(1)$ &
\textbf{0.8274\,${\pm}$\,0.011} & \textbf{0.1187\,${\pm}$\,0.004} & \textbf{0.0912\,${\pm}$\,0.003} \\
\midrule

\multicolumn{6}{@{}l@{}}{\textit{SADT}} \\
EEGNet        & CNN           & $O(1)$ & 0.5983\,${\pm}$\,0.024 & 0.2138\,${\pm}$\,0.011 & 0.1689\,${\pm}$\,0.009 \\
TSception     & CNN           & $O(1)$ & 0.6347\,${\pm}$\,0.021 & 0.2012\,${\pm}$\,0.010 & 0.1571\,${\pm}$\,0.008 \\
EEG-Conformer & CNN+Trf       & $O(T)$ & \underline{0.7058\,${\pm}$\,0.017} & \underline{0.1837\,${\pm}$\,0.008} & \underline{0.1425\,${\pm}$\,0.006} \\
RGNN          & GNN (fixed)   & $O(1)$ & 0.6715\,${\pm}$\,0.020 & 0.1936\,${\pm}$\,0.009 & 0.1512\,${\pm}$\,0.007 \\
NHGNet        & GNN (learned) & $O(1)$ & 0.6896\,${\pm}$\,0.018 & 0.1891\,${\pm}$\,0.009 & 0.1473\,${\pm}$\,0.007 \\
Compact GRU   & RNN           & $O(1)$ & 0.5672\,${\pm}$\,0.027 & 0.2251\,${\pm}$\,0.012 & 0.1786\,${\pm}$\,0.010 \\
\addlinespace[2pt]
\rowcolor{blue!7}
\textbf{HASTE (ours)} & Topo+Attn & $O(1)$ &
\textbf{0.7512\,${\pm}$\,0.014} & \textbf{0.1694\,${\pm}$\,0.007} & \textbf{0.1318\,${\pm}$\,0.005} \\
\midrule

\multicolumn{6}{@{}l@{}}{\textit{MPD-DF}} \\
EEGNet        & CNN           & $O(1)$ & 0.5623\,${\pm}$\,0.029 & 0.2284\,${\pm}$\,0.013 & 0.1815\,${\pm}$\,0.011 \\
TSception     & CNN           & $O(1)$ & 0.5941\,${\pm}$\,0.026 & 0.2153\,${\pm}$\,0.012 & 0.1697\,${\pm}$\,0.010 \\
EEG-Conformer & CNN+Trf       & $O(T)$ & \underline{0.6689\,${\pm}$\,0.021} & \underline{0.1971\,${\pm}$\,0.010} & \underline{0.1533\,${\pm}$\,0.008} \\
RGNN          & GNN (fixed)   & $O(1)$ & 0.6348\,${\pm}$\,0.024 & 0.2059\,${\pm}$\,0.011 & 0.1612\,${\pm}$\,0.009 \\
NHGNet        & GNN (learned) & $O(1)$ & 0.6527\,${\pm}$\,0.022 & 0.2017\,${\pm}$\,0.010 & 0.1576\,${\pm}$\,0.008 \\
Compact GRU   & RNN           & $O(1)$ & 0.5281\,${\pm}$\,0.032 & 0.2396\,${\pm}$\,0.014 & 0.1927\,${\pm}$\,0.012 \\
\addlinespace[2pt]
\rowcolor{blue!7}
\textbf{HASTE (ours)} & Topo+Attn & $O(1)$ &
\textbf{0.7143\,${\pm}$\,0.018} & \textbf{0.1826\,${\pm}$\,0.009} & \textbf{0.1423\,${\pm}$\,0.007} \\
\bottomrule
\end{tabular}
\end{table}

\begin{table}[!t]
\centering
\caption{Statistical significance of HASTE vs.\ each baseline
(two-sided Wilcoxon signed-rank test on Fisher-$z$ transformed per-subject PCC,
per-subject values averaged over three seeds).}
\label{tab:significance}
\renewcommand{\arraystretch}{1.12}
\setlength{\tabcolsep}{5pt}
\footnotesize
\begin{tabular}{lccc}
\toprule
\textbf{Baseline}
& \textbf{SEED-VIG} ($N{=}23$)
& \textbf{SADT} ($N{=}27$)
& \textbf{MPD-DF} ($N{=}50$) \\
\midrule
Compact GRU
& $<$0.001$^{***}$ & $<$0.001$^{***}$ & $<$0.001$^{***}$ \\
EEGNet
& $<$0.001$^{***}$ & $<$0.001$^{***}$ & $<$0.001$^{***}$ \\
TSception
& $<$0.001$^{***}$ & $<$0.001$^{***}$ & $<$0.001$^{***}$ \\
RGNN
& \phantom{$<$}0.002$^{**}$\phantom{0}
& $<$0.001$^{***}$ & $<$0.001$^{***}$ \\
NHGNet
& \phantom{$<$}0.004$^{**}$\phantom{0}
& \phantom{$<$}0.002$^{**}$\phantom{0}
& $<$0.001$^{***}$ \\
EEG-Conformer
& \phantom{$<$}0.012$^{*}$\phantom{00}
& \phantom{$<$}0.005$^{**}$\phantom{0}
& $<$0.001$^{***}$ \\
\bottomrule
\multicolumn{4}{l}{\scriptsize $^{***}\,p<0.001$;\quad $^{**}\,p<0.01$;\quad $^{*}\,p<0.05$.}
\end{tabular}
\end{table}

To assess statistical reliability, we conduct two-sided Wilcoxon signed-rank tests
on Fisher-$z$ transformed per-subject PCC values (HASTE vs.\ each baseline,
with per-subject PCC averaged over three seeds).
Table~\ref{tab:significance} reports the resulting $p$-values.
HASTE significantly outperforms all baselines on all three datasets ($p<0.05$ in all 18 comparisons).
For SADT, where some subjects have multiple sessions, per-subject PCC is obtained by averaging over that subject's sessions before applying the signed-rank test.
EEG-Conformer is the strongest competing baseline and therefore serves as the reference comparison highlighted in Fig.~\ref{fig:persubject}.
The smallest margin occurs on SEED-VIG against EEG-Conformer, where the difference remains statistically significant ($p=0.012$).

Fig.~\ref{fig:persubject} further visualizes the per-subject PCC distributions.
HASTE attains the highest median PCC on all three datasets, and its distributions remain relatively concentrated,
indicating more stable cross-subject performance than the weaker baselines.
The separation from competing methods is especially clear on SADT and MPD-DF,
while on SEED-VIG HASTE still preserves the highest center despite a partial overlap with EEG-Conformer.

\begin{figure}[!t]
\centering
\includegraphics[width=0.98\linewidth]{fig_persubject_boxplot.pdf}
\caption{Per-subject PCC distributions across three datasets for seven methods.
Each method is shown as a half-violin with individual subject points, median line, and mean marker.
(a)~SEED-VIG ($C{=}17$, $N{=}23$).
(b)~SADT ($C{=}30$, $N{=}27$).
(c)~MPD-DF ($C{=}32$, $N{=}50$).}
\label{fig:persubject}
\end{figure}

\subsection{Streaming Scalability}
\label{sec:exp_scalability}

Fig.~\ref{fig:scaling_T} shows two clearly different scaling regimes as session length $T$ increases from 100 to 10{,}000 windows.
HASTE, NHGNet, RGNN, TSception, EEGNet, and Compact GRU remain nearly flat,
indicating constant-state streaming behavior with respect to $T$.
In particular, HASTE stays around 3\,ms per update and around 50\,MB peak GPU memory across the full range,
consistent with its bounded-state design (Section~\ref{sec:method_objective}).

By contrast, EEG-Conformer grows sharply with $T$ due to full-history causal attention.
Its latency increases from only a few milliseconds at $T=10^2$ to well above $10^2$\,ms at $T=10^4$,
and its peak memory rises from roughly 50\,MB to about 1.6\,GB.
At long sessions, this yields a substantial efficiency gap in favor of HASTE.
These results show that HASTE combines the highest accuracy with predictable streaming latency and memory on long recordings.

\begin{figure}[!t]
\centering
\includegraphics[width=0.98\linewidth]{fig_scaling_T.pdf}
\caption{Streaming scalability for seven methods.
(a)~Per-update latency and (b)~peak GPU memory versus session length $T$ (windows).
Six methods achieve $O(1)$ scaling; EEG-Conformer exhibits $O(T)$ growth due to full-history causal attention.}
\label{fig:scaling_T}
\end{figure}

\subsection{Ablations and Additional Analyses}
\label{sec:exp_ablation}

\paragraph{Component ablation.}
All ablations are conducted on SEED-VIG under the same LOSO streaming protocol.
Table~\ref{tab:ablation} removes one component at a time.
Removing the topology constraint (dense attention) yields the largest PCC drop,
suggesting that the feasible neighborhood is the strongest inductive bias in the full model.
Fixing the prior-mixing gate at a constant value also degrades performance,
indicating that adaptive prior mixing is beneficial.
Uniform band weights reduce accuracy, supporting time-varying band relevance in streaming.
Setting $\beta{=}0$ reduces PCC even with masking retained, indicating that log-bias injection is not redundant with feasibility masking.
Replacing CLTA with full-history causal attention changes accuracy only slightly but sacrifices the favorable scaling in Fig.~\ref{fig:scaling_T},
which motivates the bounded temporal design.
Finally, removing $\mathcal{L}_{\mathrm{pred}}$ slightly worsens all regression metrics,
consistent with the role of trajectory regularization.


\begin{table}[!t]
\centering
\caption{Component ablation on SEED-VIG (LOSO streaming, mean\,$\pm$\,std over 3 seeds). Second line shows $\Delta$ relative to the full HASTE model.}
\label{tab:ablation}

\renewcommand{\arraystretch}{1.08}
\setlength{\tabcolsep}{0pt}
\footnotesize

\begin{tabular}{@{}l@{\hspace{7pt}}r@{\hspace{7pt}}r@{\hspace{7pt}}r@{}}
\toprule
\textbf{Variant} & \textbf{PCC}$\uparrow$ & \textbf{RMSE}$\downarrow$ & \textbf{MAE}$\downarrow$ \\
\midrule
\rowcolor{blue!7}
\textbf{HASTE (full)} &
\textbf{0.8274\,${\pm}$\,0.011} &
\textbf{0.1187\,${\pm}$\,0.004} &
\textbf{0.0912\,${\pm}$\,0.003} \\
\midrule

w/o Topo (dense attn) &
\abcell{0.7918\,${\pm}$\,0.014}{-0.036} &
\abcell{0.1296\,${\pm}$\,0.006}{+0.011} &
\abcell{0.0998\,${\pm}$\,0.005}{+0.009} \\

Fixed gate ($g_t\equiv 0$) &
\abcell{0.8041\,${\pm}$\,0.013}{-0.023} &
\abcell{0.1258\,${\pm}$\,0.005}{+0.007} &
\abcell{0.0968\,${\pm}$\,0.004}{+0.006} \\

w/o Fuse &
\abcell{0.8089\,${\pm}$\,0.012}{-0.019} &
\abcell{0.1243\,${\pm}$\,0.005}{+0.006} &
\abcell{0.0956\,${\pm}$\,0.004}{+0.004} \\

w/o Bias ($\beta{=}0$) &
\abcell{0.8112\,${\pm}$\,0.012}{-0.016} &
\abcell{0.1234\,${\pm}$\,0.005}{+0.005} &
\abcell{0.0949\,${\pm}$\,0.004}{+0.004} \\

w/o CLTA (full-history) &
\abcell{0.8248\,${\pm}$\,0.011}{-0.003} &
\abcell{0.1194\,${\pm}$\,0.004}{+0.001} &
\abcell{0.0919\,${\pm}$\,0.004}{+0.001} \\

w/o $\mathcal{L}_{\mathrm{pred}}$ &
\abcell{0.8207\,${\pm}$\,0.012}{-0.007} &
\abcell{0.1209\,${\pm}$\,0.005}{+0.002} &
\abcell{0.0931\,${\pm}$\,0.004}{+0.002} \\

\bottomrule
\end{tabular}
\end{table}
\paragraph{Sensitivity analysis.}
Fig.~\ref{fig:sensitivity}(a) shows that PCC improves rapidly as the look-back length $L$ increases from 3 to 10,
and then largely plateaus for larger windows.
Meanwhile, per-update latency grows approximately linearly with $L$, from about 1.5\,ms to about 5\,ms.

Fig.~\ref{fig:sensitivity}(b) shows that performance peaks around $k=5$
and remains relatively stable in the plateau region $k\in[4,7]$,
while both overly sparse and overly dense neighborhoods reduce accuracy.

\begin{figure}[!t]
\centering
\includegraphics[width=0.98\linewidth]{fig_sensitivity_L_k.pdf}
\caption{Sensitivity analysis on SEED-VIG.
(a)~PCC versus look-back length $L$ (bars) with per-update latency (line, right axis).
(b)~PCC versus physical $k$NN degree $k$.}
\label{fig:sensitivity}
\end{figure}

\paragraph{Prior strength and prior--attention alignment.}
To empirically validate the tunability implied by the log-bias form (Eq.~\eqref{eq:monotonic}),
we fix $\beta$ to several constants (bypassing the softplus parameterization) and measure both regression performance
and \emph{prior--attention alignment}.
We quantify alignment via the average Spearman rank correlation between the prior and attention rows restricted to the feasible neighborhood:
\begin{equation}
\rho_{\mathrm{align}}
= \frac{1}{TC}\sum_{t=1}^{T}\sum_{i=1}^{C}
  \mathrm{Spearman}\!\bigl(\mathbf{p}_{i,t},\,
  \boldsymbol{\pi}_{i,t}\bigr),
\quad
\mathbf{p}_{i,t}
\triangleq \mathbf{P}_t\bigl(i,\mathcal{N}(i)\bigr).
\label{eq:align_metric}
\end{equation}

Table~\ref{tab:beta_sweep} reports results on SEED-VIG.
As $\beta$ increases, $\rho_{\mathrm{align}}$ increases monotonically, confirming that the additive log-bias provides a continuously tunable knob
that controls how strongly spatial attention follows the knowledge prior.
Performance peaks at an intermediate bias strength, suggesting that overly strong priors can over-constrain data-driven adaptation.

\begin{table}[!t]
\centering
\caption{Prior strength sweep on SEED-VIG (LOSO streaming, mean\,$\pm$\,std over 3 seeds).
$\rho_{\mathrm{align}}$: Spearman rank correlation between the fused knowledge prior and spatial attention (Eq.~\ref{eq:align_metric}).
Learned denotes the full HASTE model with softplus-parameterized $\beta$.}
\label{tab:beta_sweep}
\renewcommand{\arraystretch}{1.15}
\setlength{\tabcolsep}{4pt}
\footnotesize
\begin{tabular}{lcccc}
\toprule
$\beta$ \textbf{setting}
& \textbf{PCC}$\uparrow$
& \textbf{RMSE}$\downarrow$
& \textbf{MAE}$\downarrow$
& $\boldsymbol{\rho}_{\mathrm{align}}$$\uparrow$ \\
\midrule
0 (no bias)
& 0.8112\,$\pm$\,.012
& 0.1234\,$\pm$\,.005
& 0.0949\,$\pm$\,.004
& 0.312\,$\pm$\,.028 \\
0.25
& 0.8165\,$\pm$\,.012
& 0.1218\,$\pm$\,.005
& 0.0938\,$\pm$\,.004
& 0.423\,$\pm$\,.025 \\
0.5
& 0.8219\,$\pm$\,.011
& 0.1202\,$\pm$\,.005
& 0.0925\,$\pm$\,.004
& 0.537\,$\pm$\,.023 \\
1.0
& 0.8261\,$\pm$\,.011
& 0.1192\,$\pm$\,.004
& 0.0916\,$\pm$\,.003
& 0.651\,$\pm$\,.021 \\
2.0
& 0.8247\,$\pm$\,.011
& 0.1196\,$\pm$\,.004
& 0.0920\,$\pm$\,.004
& 0.762\,$\pm$\,.019 \\
4.0
& 0.8158\,$\pm$\,.013
& 0.1224\,$\pm$\,.005
& 0.0941\,$\pm$\,.004
& 0.876\,$\pm$\,.016 \\
8.0
& 0.8023\,$\pm$\,.014
& 0.1267\,$\pm$\,.006
& 0.0973\,$\pm$\,.005
& 0.934\,$\pm$\,.012 \\
\midrule
Learned (softplus)
& \textbf{0.8274\,$\pm$\,.011}
& \textbf{0.1187\,$\pm$\,.004}
& \textbf{0.0912\,$\pm$\,.003}
& 0.683\,$\pm$\,.020 \\
\bottomrule
\end{tabular}
\end{table}

\subsection{Case Studies}
\label{sec:exp_case}

Fig.~\ref{fig:case_curve} shows a representative long session from SEED-VIG.
For visual clarity, the window-level predictions (originally at 1\,Hz due to 1\,s windows) are subsampled by a factor of 8
to match the 8\,s label cadence.

Panel~(a) shows that the HASTE prediction closely follows the overall rise of the ground-truth PERCLOS trajectory
while remaining smooth and tracking local fluctuations.

Panel~(b) reveals a gradual redistribution of band relevance over the session:
the beta and alpha bands occupy the largest proportions in the earlier stage,
whereas the delta and theta bands increase later and beta-band relevance decreases substantially;
alpha-band relevance remains consistently prominent throughout the session.

Panel~(c) shows that the prior-mixing gate starts at a higher level and gradually decreases over time.
In the plotted trajectory, the earlier part of the session is more anatomy-dominant,
whereas the later part becomes more functional-dominant.

\begin{figure*}[!t]
\centering
\includegraphics[width=0.98\linewidth]{fig_case_curve_alpha_gate.pdf}
\caption{Case study on SEED-VIG: streaming prediction trajectory over a representative session.
(a)~Ground-truth PERCLOS and HASTE prediction.
(b)~Band relevance weights $\boldsymbol{\alpha}_t$ (stacked area).
(c)~Prior-mixing gate $g_t$; in the figure as shown, larger values correspond to stronger anatomical contribution,
whereas smaller values correspond to stronger functional contribution.}
\label{fig:case_curve}
\end{figure*}

Fig.~\ref{fig:heatmaps} visualizes stage-dependent band-specific coupling patterns on SEED-VIG.
In the alert and transition stages, the beta and gamma bands exhibit the most pronounced posterior/occipital hotspots,
with the alpha band showing a weaker but still visible posterior concentration.

In the fatigue stage, the spatial pattern shifts:
the delta and theta bands display pronounced bilateral frontal enhancement,
while the alpha band retains a localized posterior hotspot and the beta/gamma patterns become less spatially dominant than in earlier stages,
although posterior concentration remains visible.

Overall, the figure suggests a redistribution from higher-frequency posterior-dominant coupling in the alert/transition stages
toward stronger low-frequency frontal coupling in fatigue.

\begin{figure*}[htbp]
\centering
\includegraphics[width=0.98\linewidth]{fig_prior_attention_heatmaps_2.pdf}
\caption{Band-specific coupling-weight patterns at three fatigue stages on SEED-VIG.
Columns correspond to the five frequency bands ($\delta,\theta,\alpha,\beta,\gamma$);
rows correspond to three fatigue stages (Alert, Transition, Fatigue).
For each stage, the upper panel shows the band-wise coupling/attention matrix,
and the lower panel shows the corresponding scalp projection of aggregated coupling intensity.
All heatmaps share the same electrode ordering and a unified color scale.}
\label{fig:heatmaps}
\end{figure*}
\section{Conclusion}
\label{sec:conclusion}

We presented HASTE, a hybrid-attention streaming topology-aware estimator for subject-independent EEG fatigue monitoring with continuous or continuously mapped labels. HASTE combines causal spectral fusion, topology-constrained adaptive spatial coupling with drift regularization, and prior-biased graph-restricted attention. It further employs causal local temporal attention with a fixed look-back window and a key--value cache to enable low-latency, step-wise inference with constant memory with respect to session length. Experiments under a unified LOSO streaming protocol on SEED-VIG, SADT, and MPD-DF show consistent gains over representative CNN, RNN, Transformer, and GNN baselines while preserving predictable streaming latency and memory on long sessions.

This work has several limitations. First, we focus on window-level DE features; extending HASTE to raw-EEG encoders and fully causal front-end feature extraction is an important next step. Second, we evaluate subject-independent LOSO generalization without explicit domain adaptation; integrating lightweight adaptation while retaining bounded streaming state is a promising direction. Finally, while the physical $k$NN graph provides a stable anatomy-feasible inductive bias, richer priors (e.g., montage-specific long-range links) may further improve performance under dense montages.
\newpage
\bibliographystyle{elsarticle-num} 
\bibliography{reference}
\end{document}
